% \lstset{
%     basicstyle=\ttfamily\small,
%     keywordstyle=\color{blue},
%     commentstyle=\color{gray},
%     breaklines=true,
%     frame=single,
%     showstringspaces=false
% }
% \definecolor{vscodeblue}{RGB}{86,156,214}
% \definecolor{vscodegreen}{RGB}{106,153,85}
% \definecolor{vscodered}{RGB}{214,157,133}
% \definecolor{vscodegray}{RGB}{128,128,128}

% \lstdefinestyle{vscode}{
% language=C,
% basicstyle=\ttfamily\small,
% keywordstyle=\color{vscodeblue},
% stringstyle=\color{vscodered},
% commentstyle=\color{vscodegreen},
% numbers=left,
% numberstyle=\scriptsize\color{vscodegray},
% stepnumber=1,
% numbersep=10pt,
% showstringspaces=false,
% showspaces=false,
% showtabs=false,
% breaklines=true,
% breakatwhitespace=false,
% tabsize=4,
% keepspaces=true,
% frame=none,
% columns=fullflexible,
% xleftmargin=12pt,
% xrightmargin=12pt,
% framexleftmargin=10pt,
% framexrightmargin=10pt,
% framextopmargin=5pt,
% framexbottommargin=5pt,
% aboveskip=1em,
% belowskip=1em
% }
% % 定义 Bash 样式
% \lstdefinestyle{Bash}{
%     language=bash,                % 设置语言
%     basicstyle=\ttfamily\small,   % 基础字体样式
%     keywordstyle=\color{blue},    % 关键字颜色
%     commentstyle=\color{gray},    % 注释颜色
%     stringstyle=\color{red},      % 字符串颜色
%     numbers=left,                 % 行号在左
%     numberstyle=\tiny\color{gray},% 行号样式
%     frame=single,                 % 边框
%     breaklines=true,              % 自动换行
%     showstringspaces=false        % 不显示字符串中的空格
% }
% !TeX program = xelatex
% lstlisting.tex — 代码块配置（增强 C 与 RISC‑V 汇编支持）
% 依赖：listings 宏包，建议配合 xcolor 使用。

% ================== 颜色定义 ==================
% 与 format.tex 中的 macblue 等协调；若已定义则跳过
\providecolor{macblue}{RGB}{0, 122, 255}
\providecolor{bglight}{RGB}{245, 245, 247}
\providecolor{darkgray}{RGB}{51, 51, 51}
\providecolor{warningred}{RGB}{255, 59, 48}

% 高亮专用颜色
\definecolor{codekeyword}{RGB}{0,122,255}      % 关键字蓝（同 macblue）
\definecolor{codestring}{RGB}{163,21,21}       % 字符串红
\definecolor{codecomment}{RGB}{0,128,0}        % 注释绿
\definecolor{codenumber}{RGB}{128,128,128}     % 行号灰
\definecolor{codebg}{RGB}{250,250,250}         % 背景浅灰
\definecolor{coderule}{RGB}{0,122,255}         % 边框线颜色

% VSCode 风格颜色（保留）
\definecolor{vscodeblue}{RGB}{86,156,214}
\definecolor{vscodegreen}{RGB}{106,153,85}
\definecolor{vscodered}{RGB}{214,157,133}
\definecolor{vscodegray}{RGB}{128,128,128}

% ================== 原有样式（保留） ==================
\lstdefinestyle{vscode}{
    language=C,
    basicstyle=\ttfamily\small,
    keywordstyle=\color{vscodeblue},
    stringstyle=\color{vscodered},
    commentstyle=\color{vscodegreen},
    numbers=left,
    numberstyle=\scriptsize\color{vscodegray},
    stepnumber=1,
    numbersep=10pt,
    showstringspaces=false,
    showspaces=false,
    showtabs=false,
    breaklines=true,
    breakatwhitespace=false,
    tabsize=4,
    keepspaces=true,
    frame=none,
    columns=fullflexible,
    xleftmargin=12pt,
    xrightmargin=12pt,
    framexleftmargin=10pt,
    framexrightmargin=10pt,
    framextopmargin=5pt,
    framexbottommargin=5pt,
    aboveskip=1em,
    belowskip=1em,
}

\lstdefinestyle{Bash}{
    language=bash,
    basicstyle=\ttfamily\small,
    keywordstyle=\color{blue},
    commentstyle=\color{gray},
    stringstyle=\color{red},
    numbers=left,
    numberstyle=\tiny\color{gray},
    frame=single,
    breaklines=true,
    showstringspaces=false,
}

% ================== RISC‑V 汇编语言定义 ==================
\lstdefinelanguage{RISCV}{
    % 基本指令集（RV32I / RV64I）
    morekeywords={
        add, addi, sub, and, or, xor, sll, srl, sra,
        slt, sltu, slti, sltiu,
        lb, lh, lw, ld, sb, sh, sw, sd,
        lui, auipc,
        jal, jalr,
        beq, bne, blt, bge, bltu, bgeu,
        ecall, ebreak,
        fence, fence.i,
        csrrw, csrrs, csrrc, csrrwi, csrrsi, csrrci
    },
    % M 扩展（乘除）
    morekeywords=[2]{
        mul, mulh, mulhsu, mulhu,
        div, divu, rem, remu
    },
    % F/D 扩展（浮点）与 V 扩展（向量）示例
    morekeywords=[3]{
        flw, fsw, fadd.s, fsub.s, fmul.s, fdiv.s, fsqrt.s,
        fmadd.s, fmsub.s, fnmadd.s, fnmsub.s,
        vadd.vv, vadd.vx, vsub.vv, vmul.vv, vfmul.vv
    },
    % 伪指令
    morekeywords=[4]{
        li, la, mv, call, ret, j, jr, nop, tail, not, neg
    },
    keywordstyle=[1]\color{codekeyword}\bfseries,
    keywordstyle=[2]\color{codekeyword!80!black},
    keywordstyle=[3]\color{vscodeblue!90!black},
    keywordstyle=[4]\color{codekeyword!70!black},
    sensitive=true,
    morecomment=[l]{\#},
    morecomment=[l]{//},
    morestring=[b]",
    morestring=[b]',
    % 寄存器强调
    emph={
        x0, zero, ra, sp, gp, tp,
        t0, t1, t2, t3, t4, t5, t6,
        s0, s1, s2, s3, s4, s5, s6, s7, s8, s9, s10, s11,
        a0, a1, a2, a3, a4, a5, a6, a7
    },
    emphstyle=\color{vscodeblue!70!black}\mdseries,
    alsoletter={.},
}

% ================== 新样式：C 语言（mac 风格） ==================
\lstdefinestyle{cstyle}{
    language=C,
    basicstyle=\ttfamily\small\color{darkgray},
    keywordstyle=\color{codekeyword}\bfseries,
    commentstyle=\color{codecomment}\itshape,
    stringstyle=\color{codestring},
    numberstyle=\tiny\color{codenumber},
    numbers=left,
    stepnumber=1,
    numbersep=8pt,
    backgroundcolor=\color{codebg},
    frame=single,
    framerule=0.6pt,
    rulecolor=\color{coderule!40},
    framesep=6pt,
    breaklines=true,
    breakatwhitespace=false,
    showstringspaces=false,
    showspaces=false,
    showtabs=false,
    tabsize=4,
    keepspaces=true,
    columns=fullflexible,
    xleftmargin=2.2em,
    framexleftmargin=1em,
    aboveskip=1em,
    belowskip=1em,
}

% ================== 新样式：RISC‑V 汇编（mac 风格） ==================
\lstdefinestyle{asmstyle}{
    language=RISCV,
    basicstyle=\ttfamily\small\color{darkgray},
    keywordstyle=\color{codekeyword}\bfseries,
    commentstyle=\color{codecomment}\itshape,
    stringstyle=\color{codestring},
    numberstyle=\tiny\color{codenumber},
    numbers=left,
    stepnumber=1,
    numbersep=8pt,
    backgroundcolor=\color{codebg},
    frame=single,
    framerule=0.6pt,
    rulecolor=\color{coderule!40},
    framesep=6pt,
    breaklines=true,
    breakatwhitespace=false,
    showstringspaces=false,
    showspaces=false,
    showtabs=false,
    tabsize=8,
    keepspaces=true,
    columns=fullflexible,
    xleftmargin=2.2em,
    framexleftmargin=1em,
    aboveskip=1em,
    belowskip=1em,
}

% ================== 便捷命令与环境 ==================
% 插入外部 C 文件（可选样式参数）
\newcommand{\cfile}[2][]{\lstinputlisting[style=cstyle,#1]{#2}}

% 插入外部 RISC‑V 汇编文件（可选样式参数）
\newcommand{\asmfile}[2][]{\lstinputlisting[style=asmstyle,#1]{#2}}

% 内嵌 C 代码块环境
\lstnewenvironment{ccode}[1][]
  {\lstset{style=cstyle,#1}}
  {}

% 内嵌汇编代码块环境
\lstnewenvironment{asmcode}[1][]
  {\lstset{style=asmstyle,#1}}
  {}

% ================== 默认全局设置 ==================
% 可根据需要注释或覆盖；这里将默认风格设为 cstyle
\lstset{
    basicstyle=\ttfamily\small,
    frame=single,
    breaklines=true,
    showstringspaces=false,
    style=cstyle
}